\section{Selección y Definición del Tema de Investigación}

El presenta trabajo de investigación sitúa dentro del área de \textbf{Arquitectura de software} y la \textbf{Ingeniería de Software empírica}, enfoncándose específicamente en los campos de calidad y mantenibilidad en sistemas basados en Inteligencia Artificial Generativa (GenAI) y Retrieved-Augmented Generation (RAG)

\subsection{Punto de Partida e Interés Temático}
Para la selección del tema, se parte desde interrogantes sobre la viabilidad y la relevancia de estudiar la degradación arquitectónica en aplicaciones de IA. Ante la creciente adopción de estos nuevos campos por parte de la industria, surgió la interrogante sobre si los errores en estas aplicaciones provienen exclusivamente del modelo o de la infraestructura de software. La revisión preliminar confirmó la existencia de información en repositorios de código abierto y literatura reciente \cite{shao_are_2024}, evidenciando que los resultados de esta investigación pueden aportar un marco de auditoria proactivo no presente en la actualidad.  

\subsection{Observación Directa y Consulta Bibliográfica}
A partir de la documentación en la literatura científica relacionada, se identificó el objeto de investigación: los defectos estructurales y anti-patrones de integración en aplicaciones RAG. Esta documentación permitió contrastar que, si bien existen taxonomías de fallos, la comunidad científica los aborda de manera reactiva \cite{shao_are_2024}.

\subsection{Delimitación y Definición del Tema}
Con el apoyo de la revisión bibliográfica y las directrices metodológicas del grupo de investigación MIS se delimitó el objeto de conocimiento. El tema central se define como la caracterización y detección estática de anti-patrones arquitectónicos en aplicaciones basadas en IA generativa y RAG, estableciendo el titulo preliminar del anteproyecto: \textbf{Caracterización y Detección de anti-patrones arquitectónicos en aplicaciones basadas en IA Generativa y RAG: Un estudio empírico}